\chapter[El Peso de la Indolencia]{El Peso de la Indolencia\\ \large La pereza de hoy labra un lecho de espinas para el mañana.}

\elegantimage{18_pereza_laboral/web/assets/hero.jpg}

\lettrine[lines=3, nindent=0.5em, findent=0.2em]{E}{}l mundo se sostiene sobre el pulso invisible y armónico de todos los seres que se esmeran en brindar su talento al servicio del prójimo. Escabullirse por sistema, defraudar la labor con indolencia y elegir holgazanear mientras los hermanos sostienen el peso, rompe el equilibrio...

Esta patética artimaña de falsear el propio sudor agota el flujo de la abundancia reservada. Todo lo que eludimos forzosamente acaba amontonándose como una losa frente a nuestras oportunidades.

Por lo tanto, quien no contribuye hoy ennobleciendo el trabajo, se topará invariablemente con que la vida le cierra todo acceso, siendo arrojado posteriormente a los pantanos cruentos de la miseria y a sufrir perpetuamente un agobiante desempleo.

\section{El Yunque y la Pereza}
\begin{elegantquote}
\textit{Un mozo ágil fue contratado en un gran taller rebosante de encargos, al amparo de un techo seguro. Sin embargo, dedicaba sus horas a esconder sus herramientas y dormitar bajo fardos, sintiéndose astuto por no cansar su cuerpo mientras los demás trabajaban hasta el ocaso...}

\textit{Pero la astucia que se burla del esfuerzo no tarda en cobrar su ironía. Tras malgastar sus años de vigor renegando de su trabajo, el engranaje de su destino se estancó oxidado para siempre.}

\textit{Llegado un nuevo amanecer, la miseria lo encontró suplicando por una sola hora de faena dura, sentado bajo la lluvia frente a inmensos portones de hierro herméticamente cerrados; deseando, ahora más que nunca, limpiar el sudor de su frente que antaño tanto despreció.}

\end{elegantquote}

\elegantimage{18_pereza_laboral/web/assets/art.jpg}

\vspace{1em}
\begin{center}
    \textbf{\Large El tiempo robado a la labor digna oxida irremediablemente el porvenir.}
\end{center}
\vspace{1em}

\section{Análisis Kármico y Traducción}
\begin{wrapfigure}{r}{0.5\textwidth}
  \vspace{-1.5em}
  \begin{center}
  \inlineelegantimage[0.45\textwidth]{18_pereza_laboral/web/assets/pasaje_original.jpg}
  \end{center}
  \vspace{-1.5em}
\end{wrapfigure}
El progreso de toda nuestra red infinita en la que existimos demanda esfuerzo enfocado humano. Omitir, engañar y estancarse pasiva o negligentemente a la hora de proporcionar nuestro valor equitativo no lastima a nuestros superiores, drena la balanza del sustento personal divino; resultando posteriormente despojado rigurosamente de todos los medios honorables de labor frente al agobiante cierre perpetuo de puertas hasta transmutar tal deshonra apática en servicio y esfuerzo vital.

